% Auto-extracted during the research cleanup pass. Original source is preserved unchanged.
% Source: integration.tex, lines 2213--end. Preserved block-outer/RM history and BCH spectra exploration; not included in the main spine.

Thus a local code with a published or certified spectrum drops into exactly the same first-moment interface
\(\sum_h A_h^{\mathrm{out}}p_h^{\mathrm{in}}(\delta)\).  The asymptotic theorem still uses the random banded outer; the
block code is best viewed as a finite-\(n\) implementation variant whose proof burden is shifted to a short local
spectrum certificate.

As a first citable test case we used the fourth-order binary Reed--Muller code
\[
RM(4,9):\qquad [512,256,32].
\]
Its exact weight distribution was tabulated by Markov and Borissov in 2025.  Equivalently, because \(RM(4,9)\) is a
Type-II self-dual binary code, its homogeneous weight enumerator has the Gleason form
\[
W_{RM}(x,y)
=
\sum_{j=0}^{21} c_j\,
\Phi(x,y)^{64-3j}\Psi(x,y)^j,
\]
where
\[
\Phi=x^8+14x^4y^4+y^8,
\qquad
\Psi=x^4y^4(x^4-y^4)^4.
\]
The integer coefficient vector \((c_0,\ldots,c_{21})\) is
\[
\begin{gathered}
1,\ -896,\ 366400,\ -90591488,\ 15125312416,\ -1804090983680,\\
158633285104000,\ -10462864561216000,\ 521844917396785920,\\
-19703740374727094272,\ 560488125281758654464,\ -11885355149550800420864,\\
184790978198497210204160,\ -2057511108351137030602752,\\
15886111690015316194099200,\ -81428598375642719259197440,\\
261053953183414744116101120,\ -481398485705822986116792320,\\
451352262788439851856297984,\ -176727086611150405833326592,\\
20139935110366512638066688,\ -279475807551445268430848 .
\end{gathered}
\]
For the first-moment evaluator, \(W_{\mathrm{loc}}(z)=W_{RM}(1,z)\).  Using this spectrum with \(k=2^{20}\),
\(M=4096\), \(N=2^{21}\), fixed-tap inner memory \(\sigma=32\), and the same global episode-cover inner bound, gives
the following low-weight diagnostics:
\[
\begin{array}{c|rrrr}
\delta & .106 & .104 & .092 & .090\\
\hline
-\log_2 \mu_{\le220}^{RM(4,9)} & 32.686 & 33.603 & 39.474 & 40.521\\
h_* & 32 & 32 & 32 & 32
\end{array}
\]
The peak is the minimum-weight local block term.  Thus \(RM(4,9)\) does not by itself give a \(40\)-bit margin at
\(\delta=.106\), but it does give a clean, fully spectrum-based checkpoint around \(\delta\approx .09\) at
\(\sigma=32\).  The important finite-\(k\) headline is
\[
k=2^{20},\qquad N=2^{21},\qquad \sigma=32,\qquad
\Pr[d_{\min}\le188743]\le2^{-40.5},
\]
or equivalently \(d_{\min}\ge188744\) with probability at least \(1-2^{-40.5}\).  In relative terms this certifies
\[
\frac{d_{\min}}{N}\ge\frac{188744}{2097152}>0.09000015
\]
except with probability at most \(2^{-40.5}\).  The nearby \(\delta=.092,.104,.106\) entries in the table are prefix
diagnostics, not full manifest certificates.  This is already much better structured than rejection sampling away the
first few random-banded outer weights: the entire finite outer spectrum is explicit.

At \(\delta=.09\) we also extended the same computation to \(h\le500\).  The total remained
\[
\log_2\mu_{\le500}^{RM(4,9)}\le -40.520728,
\]
again with peak \(h=32\).  On the computed tail window \(300\le h\le500\), the worst adjacent log-ratio was
\(-0.990724\).  More importantly, the analytic bounded-\(r\) low-slot ratio below gives a weaker but still sufficient
post-prefix geometric residual below \(2^{-500.579}\).  The remaining certificate task for the RM finite-\(n\)
checkpoint is therefore quite focused: keep the product RM spectrum in the residual ledger and use the same
post-prefix domination split developed above for the random-banded outer.

The first residual constants are already favorable.  Using the RM product generating bound at \(z=.4\),
\[
\log_2(W_{\mathrm{out}}(.4)-1)=33.494442,
\]
the bounded-\(r\) low-slot ratio at the handoff \(H=500\), \(R=12\), and \(\theta=.35\) is
\[
\log_2 q_{\mathrm{low}}\le -0.119386.
\]
The corresponding high-slot exponent is \(0.062159\), while the large-\(r\) low-slot ledger is maximized at
\((h,r,c)=(500,13,26)\) and equals \(-271.720148\), or \(-186.135185\) after the coarse \(3N^4\) union factor.  The
linear-window handoff from \(H/N\) to \(\eta_{\mathrm{crit}}\) still has a worst gap below \(-0.00290\).  Thus the
RM checkpoint appears to have a full finite-certificate path, not merely a plotted finite prefix.

\begin{theorem}[Finite RM-block checkpoint at relative distance \(0.09\)]
\label{thm:finite-rm-block-checkpoint-009}
Let the outer code be the direct sum of \(4096\) copies of \(RM(4,9)\), viewed as a binary
\([512,256,32]\) block code.  Thus the outer message length is \(k=4096\cdot256=2^{20}\) and the
outer output length is
\[
N=4096\cdot512=2^{21}.
\]
Concatenate this outer code through a uniform interleaver with the fixed-tap dense recursive inner ensemble with
memory \(m=32\), and set
\[
d=\lfloor .09\,N\rfloor=188743.
\]
The finite checkpoint certificate gives
\[
\Pr[d_{\min}(\Csc)\le d]
\le
\mu_{\mathrm{cert}}^{RM}
\le
2^{-40.5}.
\]
Equivalently, with probability at least \(1-2^{-40.5}\) over the interleaver and inner ensemble,
\[
d_{\min}(\Csc)\ge188744,
\]
so the relative minimum distance is larger than \(0.09\).
\end{theorem}

\begin{proof}
Apply Lemma~\ref{lem:block-outer-first-moment-interface} with \(C_{\mathrm{loc}}=RM(4,9)\).  Let
\[
W_{RM}(z)=\sum_j A_j^{RM}z^j
\]
be the exact weight enumerator of \(RM(4,9)\).  The Markov--Borissov table gives this local spectrum, and the committed
copy is checked to sum to \(2^{256}\), have maximum weight \(512\), and have minimum nonzero weight \(32\).  Therefore
the direct-sum outer has generating function \(W_{RM}(z)^{4096}\), with the all-zero word removed.

The finite prefix \(h\le500\) uses the exact local RM spectrum, the product generating-function outer bound, and the
global episode-cover inner bound with exact final-survivor accounting through the computed prefix.  The checkpoint
artifact gives
\[
\log_2 \sum_{1\le h\le500} A_h^{\mathrm{out}}p_h^{\mathrm{in}}(.09)
\le
-40.520728.
\]
For the residual after \(H=500\), apply Lemma~\ref{lem:finite-block-certificate-template} with \(R=12\) and
\(\theta=.35\).  The bounded-\(r\) low-slot, high-slot, large-\(r\) low-slot, and linear-window handoff reductions are
the same as in Lemma~\ref{lem:checkpoint-residual-split-0106}, but with the fixed-tap banded outer generating function
replaced by \(W_{RM}(z)^{4096}-1\) at \(z=.4\).  The residual verifier checks the displayed constants:
\[
\log_2 q_{\mathrm{low}}\le -0.119386,\qquad
E_{\mathrm{hi}}\ge0.062159,
\]
large-\(r\) low-slot union below \(2^{-186.135185}\), and linear handoff gap below \(-0.00290\).  The finite geometric
tail certificate records the observed post-prefix residual below \(2^{-504.095}\), while the weaker analytic
low-slot ratio already gives a post-prefix residual below \(2^{-500.579}\).  Combining the finite prefix and all residual
pieces keeps the certified first moment below \(2^{-40.5}\).  Markov's inequality gives the stated failure-probability
bound.
\end{proof}

This also explains what a stronger BCH-like block would need to provide.  A distance-only certificate is not enough:
the pessimistic bound \(W_{\mathrm{loc}}(z)\le 1+(2^b-1)z^{d_0}\) is far too loose for \(b=256\).  What is needed is a
published or certified local spectrum envelope.  Primitive binary BCH arithmetic gives a natural candidate near
length \(511\): the narrow-sense code with designed distance \(61\) has dimension \(259\) by the cyclotomic-coset
degree count, so a \(256\)-dimensional subcode with distance at least \(61\) is plausible.  Such a subcode would use
\[
4096\text{ blocks},\qquad k=2^{20},\qquad N=4096\cdot511=2093056.
\]
The random-like local-spectrum diagnostic for a \([511,256,61]\)-shaped block, at \(\sigma=32\) and \(\delta=.106\),
gives
\[
-\log_2\mu_{\le220}^{\mathrm{randlike}}\ge 96.675835,
\qquad h_*=89,
\]
with \(d=\lfloor .106N\rfloor=221863\).  This is not a certificate, but it says the BCH direction has enough apparent
headroom.  A uniform multiplicative loss of about \(2^{56.75}\) on the nonzero random-like local spectrum would still
leave a \(40\)-bit margin in this prefix diagnostic.  The remaining missing object is therefore quite specific: an
upper spectrum envelope for the
\([511,259,\ge61]\) BCH parent, or directly for a certified \([511,256,\ge61]\) subcode.  The parent spectrum would
already upper-bound any subcode spectrum, losing at most the three-dimensional overcount at the total-codeword level.
The known asymptotic BCH-spectrum theorems are useful orientation, but not yet a finite certificate at this length:
plugging the explicit error factor from the Krasikov--Litsyn BCH spectrum bound into the native \(n=511\),
designed-distance-\(61\) candidate gives error allowances of roughly \(2^{2258}\) at \(h=61\) and \(2^{2969}\) at
\(h=220\) when the denominator parameter is matched to the generator degree.  Thus the current BCH task is not to cite
an asymptotic random-like spectrum theorem verbatim; it is to obtain a short finite envelope for the relevant
low-weight parent spectrum, or to choose a nearby local code whose finite spectrum is already tabulated.
The first natural adaptation of the Sidelnikov/Krasikov--Litsyn proof is a dual-window MacWilliams bound: if the
nonzero dual weights lie in a central window \(W\), then
\[
A_h(C)\le 2^{-r}\binom{511}{h}+\max_{j\in W}|K_h(j)|.
\]
This support-only version is much too weak in our regime.  The Carlitz--Uchiyama half-width for \(t\approx30\) already
covers the entire length, and even the unrealistic window \(W=\{256\}\) would spend about \(116.18\) bits of spectrum
slack at \(h=61\) and \(80.78\) bits at \(h=89\), exceeding the \(56.75\)-bit budget for the
\([511,259,\ge61]\) target.  Thus the adaptable part of the old BCH spectra proof is not merely dual concentration; it
must be the moment/cancellation step that uses the BCH zero set to make the nonzero dual contribution cancel across
weights.

The finite target can be stated cleanly in MacWilliams form.  Let \(C\) be the \([511,259]\) BCH parent, let
\(r=252\), and write
\[
A_h(C)=2^{-r}\left(\binom{511}{h}+S_h\right),
\qquad
S_h=\sum_{j>0} A_j(C^\perp)K_h(j),
\]
where \(K_h\) is the binary Krawtchouk polynomial.  Since the actual implementation target is a
\([511,256]\) subcode, using the parent spectrum already spends three local bits.  The \(56.75\)-bit random-like slack
budget is therefore met by the parent if, for all local weights used by the prefix certificate,
\[
\log_2\left(1+\frac{S_h}{\binom{511}{h}}\right)\le 53.75,
\]
or more pessimistically if
\[
\log_2\left(1+\frac{|S_h|}{\binom{511}{h}}\right)\le 53.75.
\]
This is the concrete cancellation theorem we need from the BCH zero set.  The lower-rate native neighbor has no
three-bit parent/subcode tax, so the analogous target for the \([511,250,63]\) code is \(62.71\) bits.

There is an equivalent subset-count formulation which may be the more useful proof object.  Let \(q=512\), identify the
\(511\) cyclic positions with \(\F_q^\times\), and let \(B(\Delta)\) be the primitive narrow-sense BCH code whose zero
set contains \(\alpha,\alpha^2,\ldots,\alpha^{\Delta-1}\).  Then
\[
A_h(B(\Delta))
=
\#\left\{
T\subseteq \F_q^\times:\ |T|=h,\quad
\sum_{x\in T}x^i=0\ \text{for every }1\le i<\Delta
\right\}.
\]
Indeed, the binary word with support \(T\) has codeword polynomial \(c(z)=\sum_{x=\alpha^\ell\in T}z^\ell\), and the
BCH zero conditions are exactly \(c(\alpha^i)=\sum_{x\in T}x^i=0\).  Thus the missing BCH theorem can also be stated as
a finite equidistribution bound for \(h\)-subsets of \(\F_{512}^\times\) under the first \(\Delta-1\) power sums.  For
\(\Delta=61\), the target is that this constrained subset count is at most \(2^{53.75}\) times the random-like
\([511,259]\) count, uniformly over the local prefix weights used by the block certificate.  This is the same content
as the MacWilliams correction bound, but it exposes the algebraic equations that a direct finite proof would need to
exploit.
In characteristic two the even constraints are Frobenius consequences of the odd ones:
\[
\left(\sum_{x\in T}x^i\right)^2=\sum_{x\in T}x^{2i}.
\]
For the \(\Delta=61\) parent, the cyclotomic-coset representatives below \(61\) are
\[
1,3,5,\ldots,31,\quad 35,37,39,41,43,45,47,\quad 51,53,55,57,59,
\]
twenty-eight size-\(9\) cosets, giving redundancy \(252\).  The \(\Delta=63\), \(k=250\) neighbor adds the coset of
\(61\), giving redundancy \(261\).  Thus a finite proof only has to control the joint distribution of these
\(28\) or \(29\) trace/power-sum constraints, not sixty unrelated equations.

The additive-character version is the bridge back to MacWilliams.  For a set \(E\) of coset representatives, let
\[
P_{\mathbf a}(x)=\sum_{e\in E}a_e x^e,\qquad a_e\in\F_{512},
\]
and let \(\psi(u)=(-1)^{\mathrm{Tr}_{\F_{512}/\F_2}(u)}\).  Orthogonality gives
\[
A_h(B(\Delta))
=2^{-9|E|}
\sum_{\mathbf a\in\F_{512}^{E}}
[z^h]\prod_{x\in\F_{512}^\times}\left(1+z\,\psi(P_{\mathbf a}(x))\right).
\]
The zero choice of \(\mathbf a\) is the random-like main term.  For \(\mathbf a\ne0\), if
\[
w(\mathbf a)=\#\{x\in\F_{512}^\times:\ \mathrm{Tr}(P_{\mathbf a}(x))=1\},
\]
then the coefficient is
\[
[z^h](1+z)^{511-w(\mathbf a)}(1-z)^{w(\mathbf a)}=K_h(w(\mathbf a)).
\]
Thus the missing cancellation theorem can be attacked either as a statement about the weight distribution of the dual
trace-code family \(\mathrm{Tr}(P_{\mathbf a}(x))\), or directly as a high-degree elementary-symmetric cancellation
over the signs \(\psi(P_{\mathbf a}(x))\).

There is also special structure at the minimum-weight edge.  If \(T\subseteq\F_{512}^\times\) has size \(h\) and
\(\sum_{x\in T}x^i=0\) for \(1\le i<h\), then Newton's identities imply, in characteristic two, that all odd elementary
symmetric functions below degree \(h\) vanish.  For odd \(h\), the monic polynomial with roots \(T\) therefore has the
form
\[
f_T(X)=X^h+e_2X^{h-2}+e_4X^{h-4}+\cdots+e_{h-1}X+e_h.
\]
Moreover \(f_T(X)=X f_T'(X)+e_h\).  Since the roots are nonzero, every root satisfies \(x f_T'(x)=e_h\).  This does
not yet count the minimum-weight words, but it is a sharper algebraic object than a generic subset of size \(h\): the
\(h=61\) and \(h=63\) boundary slices are constrained root sets of this special polynomial family.
For the \(\Delta=61\) parent, the first free power sum at the boundary is \(p_{61}\).  Newton gives
\(p_{61}=e_{61}\), so it is nonzero for a weight-\(61\) support in \(\F_{512}^\times\).  Since \(\gcd(61,511)=1\),
scaling \(T\mapsto \beta T\) lets one normalize \(p_{61}=1\).  This is the same locator-polynomial reduction used in
the Augot--Charpin--Sendrier minimum-distance approach to BCH codes.  For the \(\Delta=63\), \(k=250\) neighbor, the
first free boundary power sum is \(p_{63}\), but \(\gcd(63,511)=7\), so scaling leaves seven nonzero classes rather
than a single normalized case.
In the normalized \(h=61\) case, writing
\[
g(Y)=Y^{30}+e_2Y^{29}+e_4Y^{28}+\cdots+e_{60},
\]
the locator polynomial is
\[
f_T(X)=Xg(X^2)+1.
\]
Thus the boundary supports are exactly the \(61\)-element subsets of \(\F_{512}^\times\) on which
\[
xg(x^2)=1
\]
for some monic degree-\(30\) polynomial \(g\).  Since \(x\mapsto x^2\) is a permutation of \(\F_{512}^\times\), any
\(31\) proposed support points determine \(g\) by interpolation; the remaining \(30\) points must then be additional
roots of the same degree-\(61\) polynomial \(Xg(X^2)+1\).  A finite boundary certificate could therefore be built by
counting, or upper-bounding, the monic degree-\(30\) polynomials \(g\) for which \(Xg(X^2)+1\) splits completely over
\(\F_{512}^\times\).  The naive \(2^{270}\) count of all such \(g\)'s is far too loose, so the splitting condition is
the essential extra structure.  Including the \(511\)-element scaling orbit, this naive boundary count is about
\(2^{278.997}\), whereas the \([511,259]\)-parent certificate would allow only about \(2^{67.131}\) weight-\(61\)
words.  Thus the splitting condition must save roughly \(212\) bits at the boundary.  The analogous naive count for
the \([511,250,63]\) neighbor is about \(215\) bits above its allowance.
The random-polynomial heuristic says this amount of saving is plausible.  A random monic degree-\(61\) polynomial over
\(\F_{512}\) splits into \(61\) distinct nonzero roots with probability
\[
\binom{511}{61}/512^{61}\approx 2^{-283.62}.
\]
Multiplying by the \(2^{270}\) normalized locator family and the \(511\)-element scaling orbit gives an expected
boundary codeword count around \(2^{-4.62}\), about \(71.75\) bits below the certificate allowance.  The analogous
\(h=63\) heuristic has about \(80.71\) bits of room.  This is only heuristic, but it suggests the boundary slice is
probably not structurally fatal; the proof needs a finite splitting-count upper bound for this special family.
Equivalently, after the change of variable \(Y=X^2\), the normalized \(h=61\) boundary words are monic degree-\(30\)
polynomials \(g(Y)\) agreeing with the word \(Y^{255}\) on \(61\) points of \(\F_{512}^\times\).  The generic
interpolation/list-size bound gives
\[
\#\{g\}\le \binom{511}{30}/\binom{61}{30},
\]
so the total boundary codeword count is at most \(2^{112.260}\) after restoring the scaling orbit.  This is rigorous
but still about \(45.13\) bits above the certificate allowance, so generic Reed--Solomon list-decoding bounds are not
yet enough.  The special received word has extra structure:
\[
Y\,(Y^{255})^2=1\quad \text{on }\F_{512}^\times.
\]
Thus the boundary problem is equivalently to count monic degree-\(30\) polynomials \(g\) for which
\[
Yg(Y)^2+1
\]
has \(61\) distinct nonzero roots.  This lacunary splitting formulation is the next finite theorem target.
One way to see the remaining gap is to use more than the \(30\) interpolation points needed to determine \(g\).  If a
valid boundary word has support size \(61\), then for any \(s\ge30\) it is counted \(\binom{61}{s}\) times by its
\(s\)-subsets.  The raw bound
\[
\binom{511}{s}/\binom{61}{s}
\]
gets worse as \(s\) grows, but an \(s\)-set is consistent with a monic degree-\(30\) interpolant only if roughly
\(s-30\) algebraic conditions over \(\F_{512}\) hold.  If those conditions behaved independently, the \(s=45\) ledger
would save \(15\cdot9=135\) bits, enough to beat the certificate allowance by about \(25\) bits.  Thus the concrete
finite problem is not to improve generic interpolation by a little; it is to prove a rank/equidistribution statement
for these overdetermined interpolation conditions.
Small exact toy cases support the random-splitting intuition.  For the analogous family \(Yg(Y)^2+1\) over
\(\F_{32}\) with \(h=5\), exact enumeration gives \(6\) completely split \(g\)'s, compared with random expectation
\(2^{2.37}\).  Over \(\F_{64}\) with \(h=7\), exact enumeration gives \(85\), compared with random expectation
\(2^{5.04}\).  These small cases are not evidence for a theorem at \(q=512\), but they do not show the kind of
large oversplitting bias that would make the boundary target implausible.
A direct \(q=512\) sampler gives the same first-order signal for overdetermined interpolation.  For random \(31\)-sets
of points in \(\F_{512}^\times\), the observed consistency rate was \(13/5000\), on the \(1/512\) scale predicted by
one independent \(\F_{512}\)-condition.  For random \(32\)-sets, \(0/5000\) successes were observed, consistent with
the much smaller \(512^{-2}\) prediction.  This is again only diagnostic, but it supports the idea that the missing
finite theorem is a rank/equidistribution statement rather than a fight against an obvious algebraic bias.
The first overdetermined condition has an explicit symmetric form.  For a \(31\)-set \(S\subset\F_{512}^\times\), the
top coefficient of the degree-\(30\) interpolation polynomial for \(Y^{255}\) on \(S\) is the divided difference
\((Y^{255})[S]\), which equals the complete homogeneous symmetric polynomial \(h_{225}(S)\).  Therefore \(S\) is
consistent with a monic degree-\(30\) \(g\) iff
\[
h_{225}(S)=1.
\]
A direct checker verifies this identity against interpolation, and a \(q=512\) sample gave \(3/1000\) successes for
this equation, again on the \(1/512\) scale.  For larger \(s\), the corresponding rank problem is that all \(31\)-subsets
of the \(s\)-set have the same divided-difference value \(1\), with only \(s-30\) independent conditions expected.
The follow-up rank-audit tool makes this interpolation formulation reproducible without enumerating the full
locator-coefficient family.  Exact \(h\)-set enumeration in the small \(m=5\) cases gives the same normalized boundary
counts as the locator enumerator:
\[
\begin{array}{c|c|c|c}
m & h & \text{consistent }h\text{-sets} & A_h\\
\hline
5 & 5 & 6 & 186\\
5 & 7 & 5 & 155 .
\end{array}
\]
For the \(m=6,\ h=13\) exact-spectrum checkpoint, randomized interpolation from \(d=6\) points recovered all
\(28\) normalized split polynomials in a \(200000\)-sample run, restoring the exact boundary count
\(28\cdot63=1764\).  This is not a deterministic proof of the \(m=6\) boundary slice, but it validates that the
rank/interpolation representation is seeing the same objects as the full BCH spectrum.  The same script samples the
first \(q=512\) overdetermined condition; a \(2000\)-sample run found \(3\) successes against the random-rank prediction
\(2000/512\).  On the next half-rate rung, \(m=7,\ h=21\), the first rank sample found \(29/5000\) successes, again on
the \(1/128\) scale.

As a check on the whole BCH-spectrum program, we started a small exact ladder.  The primitive narrow-sense spectra are
computed directly from the binary power-sum parity checks by enumerating the nullspace.  The first exact data points
are:
\[
\begin{array}{c|c|c|c|c|c}
m & n & \Delta & k_{\mathrm{BCH}} & d_{\min} & \max_{0<h<n}\log_2(A_h/A_h^{\mathrm{rand}})\\
\hline
5 & 31 & 5 & 21 & 5 & 0.165\\
5 & 31 & 7 & 16 & 7 & 1.011\\
6 & 63 & 13 & 30 & 13 & 0.591\\
6 & 63 & 15 & 24 & 15 & 1.551
\end{array}
\]
Here \(A_h^{\mathrm{rand}}=2^{k-n}\binom{n}{h}\).  These are small codes, but they are useful calibration: the exact
low-weight BCH spectra are close to random-like rather than exhibiting the kind of large inflation that would doom the
block-outer upgrade.  They also validate the power-sum/nullspace construction used by the larger \(n=511\) target.

There is also a slightly lower-rate native neighbor.  The Canteaut--Chabaud length-\(511\) table identifies the
narrow-sense BCH code with dimension \(250\) and designed distance \(63\) as having true minimum distance \(63\).  A
random-like diagnostic for a \([511,250,63]\)-shaped block gives
\[
-\log_2\mu_{\le220}^{\mathrm{randlike}}\ge102.653059,
\qquad h_*=89,
\]
at the same \(\sigma=32\), \(\delta=.106\) checkpoint, with \(4195\) local blocks and
\(N=2143645\).  Its uniform random-like slack budget before the \(40\)-bit margin disappears is about \(62.71\) bits.
This neighbor is therefore not worse from the first-moment point of view.  The tradeoff is conceptual rather than
numerical: it gives up six local message bits per block but starts from a better-documented minimum-distance line.

\paragraph{Research proposal: affine-coset spectra ladder.}
The most plausible route to accurate finite BCH spectra is to scale the method used for the exact length-\(128\) and
length-\(256\) extended primitive BCH weight distributions, rather than to enumerate primitive codewords directly.
The extended code is invariant under the affine group of \(\F_{2^m}\), while the primitive code only exposes the cyclic
subgroup.  The proposed ladder is:
\[
\begin{array}{ll}
1. & \text{Reproduce the published }[128,64]\text{ extended BCH spectrum and derive }[127,64]\text{ by puncturing.}\\
2. & \text{Implement affine-orbit/coset representatives on the }m=7\text{ extended code and validate against the table.}\\
3. & \text{Move to }m=8\text{, where published length-}256\text{ extended BCH spectra provide a harder validation rung.}\\
4. & \text{For }m=9\text{, compute a low-weight prefix envelope rather than the full distribution, using nested subcodes.}
\end{array}
\]
The computational kernel should follow the local-weight-distribution viewpoint of Yasunaga--Fujiwara: view the code as
cosets of a subcode, partition those cosets under the automorphism group, and compute the weight distribution for one
representative per orbit.  The boundary interpolation-rank tool remains useful for the minimum-weight slice, while a
dual trace/MacWilliams sampler remains useful as a prefix diagnostic; the affine-coset ladder is the primary candidate
for a reproducible spectrum table at message dimension around \(128\).
The first table-audit rung confirms the extension relation exactly on the public Desaki tables:
\([128,64,22]\) derives the published \([127,64,21]\) primitive spectrum with no mismatches.  On the \(m=8\) rung, the
public Okayama table currently links extended length-\(256\) spectra up to dimension \(63\) and from dimension \(207\)
upward, with the central dimensions \(71,\ldots,199\) listed but not linked.  The available \([256,63]\) table likewise
derives the published \([255,63]\) primitive table with no mismatches.  Thus the public database is enough to validate
the import and extension algebra, but not enough by itself for a \(k\approx128\) block; the affine-coset computation is
still needed for the central dimension.
As a direct block-outer trial, we also imported the Okayama/Desaki exact weight distribution
\(\texttt{EBCH128\_64.wd}\).  The table sums to \(2^{64}\), has length \(128\), and has \(d_{\min}=22\).  Plugging this
exact \([128,64,22]\) local spectrum into the same \(k=2^{20}\), \(N=2^{21}\), \(\sigma=32\) block-outer first-moment
prefix gives
\[
\begin{array}{c|c|c|c}
\delta & d=\lfloor \delta N\rfloor & -\log_2 \mu_{\le220} & h_*\\
\hline
.106 & 222298 & 12.382683 & 25\\
.09 & 188743 & 18.592737 & 24
\end{array}
\]
Thus this specific exact EBCH table is a useful import/sanity rung, but it is not competitive with the RM checkpoint:
its local floor \(22\) leaves the small-block-weight contribution too large for the current finite-\(n\) target.
The first affine-orbit sanity check works on exact-enumerable \(m=5\) BCH codes.  The important bookkeeping point is
that affine orbits live in the extended code: an extended weight \(W\) slice combines primitive weights \(W-1\) and
\(W\).  With this convention, the \([31,21,5]\) code has one affine orbit of \(992\) extended weight-\(6\) words, and
the \([31,16,7]\) code has five affine orbits of \(124\) extended weight-\(8\) words.

The reproducibility entry points are the block-outer probe, BCH-parameter script, and the RM Gleason/checkpoint/residual
verifiers in \(\texttt{scripts/}\).  The RM checkpoint constants are also collected in the certificate manifest
\[
\texttt{scripts/certificates/rm512\_256\_sig32\_delta009.json},
\]
which is audited by
\[
\texttt{python scripts/verify\_block\_certificate\_manifest.py}.
\]
For a new spectrum-csv block candidate, the corresponding manifest is generated from the prefix and tail artifacts by
\[
\texttt{python scripts/create\_block\_certificate\_manifest.py}.
\]
The BCH asymptotic-bound diagnostic above is recorded by
\[
\texttt{python scripts/check\_bch\_kl\_spectrum\_bound.py}.
\]
The corresponding random-like slack budget is recorded by
\[
\texttt{python scripts/check\_bch\_spectrum\_slack\_budget.py}.
\]
The dual-window MacWilliams audit is recorded by
\[
\texttt{python scripts/check\_bch\_dual\_window\_envelope.py}.
\]
The parent/subcode correction target is recorded by
\[
\texttt{python scripts/check\_bch\_macwilliams\_target.py}.
\]
The Newton boundary normalization ledger is recorded by
\[
\texttt{python scripts/check\_bch\_newton\_boundary.py}.
\]
The corresponding trivial boundary-count gap is recorded by
\[
\texttt{python scripts/check\_bch\_boundary\_trivial\_bounds.py}.
\]
The random-splitting boundary heuristic is recorded by
\[
\texttt{python scripts/check\_bch\_boundary\_split\_heuristic.py}.
\]
The finite interpolation bound is recorded by
\[
\texttt{python scripts/check\_bch\_boundary\_interpolation\_bound.py}.
\]
The higher-interpolation rank budget is recorded by
\[
\texttt{python scripts/check\_bch\_boundary\_higher\_interp\_budget.py}.
\]
Small-field exact splitting sanity checks are recorded by
\[
\texttt{python scripts/check\_bch\_boundary\_smallfield.py}.
\]
The \(q=512\) interpolation-consistency sampler is recorded by
\[
\texttt{python scripts/check\_bch\_boundary\_interp\_sampling.py}.
\]
The divided-difference form of the first consistency condition is recorded by
\[
\texttt{python scripts/check\_bch\_boundary\_divdiff.py}.
\]
The overdetermined interpolation-rank audit is recorded by
\[
\texttt{python scripts/bch\_boundary\_rank\_audit.py}.
\]
For example, the \(m=6,\ h=13\) recovery check above is reproduced by
\[
\texttt{python scripts/bch\_boundary\_rank\_audit.py --m 6 --h 13 --compare-spectrum scripts/bch63\_30\_delta13\_spectrum.csv recover-polys --trials 200000 --seed 23}.
\]
The published-table import and extended-to-primitive validation rung is recorded by
\[
\texttt{python scripts/bch\_wd\_table\_tools.py}.
\]
Single \(\texttt{.wd}\) files in the public Okayama format are converted to block-probe CSV spectra by
\[
\begin{array}{l}
\texttt{python scripts/import\_wd\_spectrum.py scripts/EBCH128\_64.wd}\\
\texttt{\quad scripts/ebch128\_64\_spectrum.csv --expect-dim 64 --expect-length 128}.
\end{array}
\]
For example, the public \(m=8\) availability audit is
\[
\texttt{python scripts/bch\_wd\_table\_tools.py --insecure-url --index-url https://isec.ec.okayama-u.ac.jp/home/kusaka/wd/ --section "Extended Binary Primitive BCH Codes" --length 256}.
\]
The first affine support-orbit sanity checks are recorded by
\[
\texttt{python scripts/bch\_affine\_orbit\_sanity.py}.
\]
The small exact BCH spectra are generated by
\[
\texttt{python scripts/exact\_bch\_spectrum\_small.py}.
\]
Their compact random-like-loss summaries are generated by
\[
\texttt{python scripts/summarize\_exact\_bch\_spectrum.py}.
\]
The first specialized enumerator is the boundary locator count.  For the two \(m=5\) exact spectra it reproduces the
minimum-weight counts without enumerating the full code:
\[
\begin{array}{c|c|c|c|c}
n & \Delta & k & \text{normalized locator hits} & A_\Delta\\
\hline
31 & 5 & 21 & 6 & 186\\
31 & 7 & 16 & 5 & 155
\end{array}
\]
where \(A_\Delta\) is obtained by restoring the \(31\)-element scaling orbit.  This validates the Newton/locator
boundary reduction against exact spectra.  The naive Python root-testing implementation is not yet fast enough for the
\(m=6\), \(\Delta=13\) boundary family, so the next tooling step is a vectorized or bitset version of the same
enumerator.
The boundary locator enumerator is
\[
\texttt{python scripts/enumerate\_bch\_boundary\_locator.py}.
\]

\paragraph{What the exact isolated statistics buy us.}
Corollary~\ref{cor:isolated-order-late} gives a rigorous next rung above the late-placement base: on the isolated
slice one can count exactly the event that the \(r\)-th isolated start lies inside the final \(L\) positions. This is
the clean theorem object behind statements like ``there is at most one genuinely early isolated impulse.'' Current
numerics show that the \(r=2\) version is still too weak by itself to certify the observed small-weight regime, but it
is now a precise proved object rather than a heuristic guess.

Corollary~\ref{cor:isolated-early-count} sharpens that viewpoint further: on the isolated slice one can count exactly
how many isolated starts fall before the final \(L\) positions. Empirically, a model of the form
\[
\sum_{j=0}^h
\PP[J_L(U)=j\mid \wt(U)=h,\ R(U)=h]\tau^j
\]
with a single \(\tau\) already tracks the exact isolated slices within about a bit on the current overlap tables.
That is encouraging because it isolates the remaining analytic burden into a per-early-episode penalty.
At the same time, a uniform safe \(\tau\) is still too loose when inserted into the large-\(k\) first moment, so the
next useful theorem likely needs a position-sensitive or \(j\)-sensitive refinement rather than a single global
constant.

The natural position-sensitive version is to assign each early isolated start position \(y\le N-L\) its own
penalty \(\tau_y(\delta)\), and then average the product of those penalties over the exact isolated-start
configuration. On the isolated slice this would take the form
\[
\frac{1}{\binom{N-h+1}{h}}
\sum_{j=0}^h
\binom{L-h+1}{h-j}\,
e_j\!\bigl(\tau_1(\delta),\dots,\tau_{N-L}(\delta)\bigr),
\]
where \(e_j\) is the \(j\)-th elementary symmetric polynomial. The uniform-\(\tau\) model is exactly the special case
\(\tau_y\equiv \tau\), so the current overlap evidence can be read as saying that the right theorem probably lives one
level above the constant-\(\tau\) simplification, not in a completely different combinatorial regime.

The most natural first choice is to take \(\tau_y(\delta)\) from the one-run bound with remaining suffix length
\(N-y+1\), namely
\[
\tau_y(\delta)
\approx
(N-y+1)2^{-m}
\;+\;
\Pr[\Binomial(N-y+1,\tfrac12)\le \delta N].
\]
That position-sensitive one-run model does explain the isolated slice better than a global constant at the very
smallest weights, but it still undershoots once \(h\) reaches the \(5\text{--}7\) window driving the current
large-\(k\) transition. So the next theorem is still not just ``plug in the one-run law per early start''; it needs
an additional multi-episode interaction term on top of that exact early-count decomposition.

\paragraph{The current best residual law.}
One especially clean candidate is suggested by Corollary~\ref{cor:isolated-early-pair-mean}: after factoring out the
position-sensitive one-run model, the remaining overlap gap is well tracked by an affine function of the exact
expected number of early-start pairs
\[
\EE\!\left[\binom{J_L(U)}{2}\,\middle|\, \wt(U)=h,\ R(U)=h\right]
=
\binom{h}{2}\frac{\binom{N-L}{2}}{\binom{N-h+1}{2}}.
\]
On the current overlap tables the residual coefficients are well fit by
\[
a(q_1,m)\approx -1.1434 + 6.1146\,q_1 - \frac{15.4212}{m},
\]
\[
b(q_1,m)\approx 0.4950 - 1.0649\,q_1 - \frac{0.7693}{m},
\]
so the residual takes the working form
\[
\text{gap bits}
\approx
a(q_1,m)
+
b(q_1,m)\,
\EE\!\left[\binom{J_L(U)}{2}\,\middle|\, \wt(U)=h,\ R(U)=h\right].
\]
This does not yet prove the right theorem, but it is a much sharper target than a free global fit: it says the next
correction likely lives on the early-start hypergeometric law itself, not on the full run count.

The new one-episode lemma confirms that \(\mathrm{Corr}_{N,h}(\delta)\) cannot be controlled by looking only at the
first early ON-episode: that route is structurally valid but numerically much too loose. The current multi-episode
span certificate is the first pessimistic correction that has the right summability shape. It is probably not sharp,
but it turns an explicit placement statistic into the proof object instead of fitting a free residual curve.

\paragraph{What must be improved next.}
The outer small-weight formula \eqref{eq:dense-dense-working-outer-small} has now largely been replaced, for
finite-\(n\) estimation purposes, by the exact low-weight systematic banded law suggested by the enumerator.
The remaining conceptual gap is therefore mostly on the inner side of the finite prefix: the current bridge residual
is summable, but the small/intermediate-weight certificate still needs the sharpest safe version of the
fixed-tap-dense recursive law.  The current best proof-oriented subproblem is to upgrade
Corollary~\ref{cor:fixedtap-weight-slice-multiepisode} by charging not only the fixed-tap terminations needed to cover
the early span, but also the fair-output cost of dense clusters of input ones that are swallowed by a live episode.
Once that is in place, the remaining task is to strengthen the exact finite-prefix computation and residual split from
the present break-even certificate to a certificate with a prescribed larger negative margin.
